% !TITLE 诗经·小雅·鹿鸣
% !AUTHOR 无名氏
% !DYNASTY 先秦
% !GENRE 四言诗
% !ORDER 2
% !SOURCE 诗经·小雅

\begin{poem}
呦呦鹿鸣，食野之苹。我有嘉宾，鼓瑟吹笙。
吹笙鼓簧，承筐是将。人之好我，示我周行。

呦呦鹿鸣，食野之蒿。我有嘉宾，德音孔昭。
视民不恌，君子是则是效。我有旨酒，嘉宾式燕以敖。

呦呦鹿鸣，食野之芩。我有嘉宾，鼓瑟鼓琴。
鼓瑟鼓琴，和乐且湛。我有旨酒，以燕乐嘉宾之心。
\end{poem}

\begin{appreciation}
《鹿鸣》是《小雅》的第一篇，也是中国最早的宴饮诗。

呦呦鹿鸣，食野之苹——鹿发现了丰美的野草，发出呦呦的鸣叫来呼唤同伴。这便是"鹿鸣"的比兴：我有嘉宾，也当如鹿之呼朋引伴，以最好的音乐和美酒来款待他们。

这首诗在后世有一个很美的延伸：鹿鸣宴。科举时代，乡试放榜后，地方官宴请新科举人，席间必奏《鹿鸣》之章，这便是"鹿鸣宴"的由来。这个传统从唐代延续到明清，千余年不曾断绝。有意思的是，宴饮的实质内容千差万别——有盛大的，也有敷衍的——但"鹿鸣"这两个字始终被保留了下来。因为它所承载的是一种朴素而珍贵的情感：人与人之间，因为宴饮聚在一起，不止是为了享乐，更是为了切磋德行。"视民不恌，君子是则是效"——宾客不只是来喝酒的，他们的品德本身就是一道风景，足以让旁人效法。

这跟今天的应酬是两回事。今天的酒局讲排场、讲利益，而《鹿鸣》里的宴饮讲的是"和乐且湛"——快乐的极致不是喧闹，而是和谐融洽。有音乐，有美酒，有彼此欣赏的人，足矣。

从《关雎》到《鹿鸣》，从爱人到友朋，《诗经》的编排自有深意。一个人先有了爱人的能力，再走出家门与朋友宴饮，这便是儒家理想中的完整人格——始于男女之情，成于君臣之义，而《鹿鸣》恰好站在这个台阶上。
\end{appreciation}
